\documentclass{article}

% typesetting
\usepackage[letterpaper, textwidth=5.5in, textheight=9.0in]{geometry}
\usepackage{setspace}
\setstretch{1.08}
\renewcommand{\paragraph}[1]{\par\medskip\noindent\textbf{#1}~---}
\frenchspacing\raggedbottom\sloppy\sloppypar

% mathematics
\usepackage{amsmath}
\newcommand{\setof}[1]{\left\{{#1}\right\}}

% collaboration
\usepackage{xcolor}
\newcommand{\hoggsay}[1]{{\textcolor{red}{Hogg say: {#1}}}}

\title{\bfseries A fully self-supervised redshift pipeline}
\author{Hogg for Chauhan}
\date{2026 October}
\begin{document}
\maketitle

\begin{abstract}\noindent
    Most data-driven redshift-determination techniques require some labeled data or ground truth.
    Here we build a fully data-driven redshift estimation system that requires precisely one---and only one---object
    in the data sample to have a known or adopted redshift.
    Everything else is learned from the properties of the data and the fundamental symmetry of the Doppler shift.
    We apply this method to the \hoggsay{DRXX} quasar catalog from the \textsl{Sloan Digital Sky Survey}.
    We find \hoggsay{what things?}
\end{abstract}

\section{Introduction}

\section{Assumptions and choices}
\paragraph{Data}
We have $N$ quasar (or galaxy) spectra $y_i$ ($1\leq i\leq N$), each of which is an ordered list $\setof{y_{ij}}_{j=1}^M$ of $M$ calibrated pixel flux or intensity measurements $y_{ij}$ on a standard Solar System barycentric rest-frame (for our purposes the ``observed'' frame) log-wavelength grid.
Each pixel $y_{ij}$ ($1\leq j\leq M$) of spectrum $y_i$, which corresponds to some observatory-frame log-wavelength $x_j \equiv \ln\lambda_j$, has an associated weight or uncertainty inverse variance $\sigma^{-2}_{ij}$.

Importantly---and as is the case for \textsl{SDSS} and \textsl{DESI} spectra, for example---we will assume that the log-wavelength grid is uniform in log wavelength.
That is, we will assume that, in these rectangular data, Doppler shifts, which are pure translations in $x$, are pure translations in the discretized log-wavelength grid.
This is an overly strong assumption; there are good methods that don't assume this.
However, since our data have this property, we are going to capitalize on it.

Fundamentally we are going to assume that these data are ``correct'' in most relevant ways:
They are correctly calibrated, they have been correctly extracted on a standardized observed-frame wavelength grid, and the uncertainty estimates are reasonable (for most pixels of most spectra).

These data are rectangular, but we won't assume that they are complete:
There might be missing or masked data, these measurements are simply given vanishing inverse uncertainty variances $\sigma^{-2}_{ij}$ such that they will have no weight in anything we do below.

\paragraph{Doppler shift}
The point of this project is to learn Doppler shifts.
As mentioned above, the Doppler shift is a pure, rigid translation in log wavelength.
That fact structures this whole project and makes the inference possible.
Each quasar $i$ will be assigned a relative Doppler shift $\Delta_i$, 
These are related to redshift $z_i$ by
\begin{align}
    \Delta_i = \ln (1 + z_i) ~.\label{eq:doppler}
\end{align}

\paragraph{Exactly one known redshift}
In order for the Doppler shifts to be identifiable, overall, we need to know the true redshift $z_i$ for exactly one object $i=i_0$.
We only need one, but we want to be very confident about that one object's redshift.

\paragraph{Low-dimensional representation}
Quasar spectra are remarkably diverse.
That said, they still live in a low-dimensional space; they can't have arbitrary spectral shape.
If they didn't live in some low-dimensional or compact space, no redshift determination would ever have been possible, by any method.
Here we will assume something even stronger, which is that there is a \emph{linear} $K$-dimensional subspace in spectral space
that gives a good representation for most quasars, or at least good enough for redshift determination.

Another way to say it is that we will assume that the expectation for the spectrum of quasar $i$ can be expressed with only $K$
amplitudes $a_{ik}$, such that
\begin{align}
    y_{ij} &= \sum_{k=1}^K a_{ik}\,g_k(x_j - \Delta_i) + s_{ij} ~,
\end{align}
where the $g_k(x)$ are rest-frame (model-frame) eigenspectra or basis functions,
they are functions of log wavelength $x$,
the $\Delta_i$ are the (relative) redshifts of the spectra $i$,
and the $s_{ij}$ are residuals or noise or individual-pixel departures from this model.

Because all spectra are on the same, uniform, observed-frame log-wavelength grid of points $x_j$,
we are going to represent these rest-frame eigenspectra $g_k(x)$ also on a uniform, rest-frame log-wavelength grid with the same grid spacing as the data.
That is, the only difference between the grid we use for the $g_k(x)$ and that used for the data is the starting and ending log wavelengths of that grid.
We discuss interpolation below.

\paragraph{Noise model}
\hoggsay{We assume that the noise has the form for which RHMF is the correct algorithm. Maybe Hilder can write this paragraph?}

\paragraph{No interpolation}
\hoggsay{We might interpolate the model expectation, but we will never interpolate the data!}

\paragraph{Initialization}
We are about to do a big optimization.
The optimization is easy at fixed Doppler shifts $\setof{\Delta_i}_{i=1}^N$.
We need to initialize the Doppler shifts to get started.
We will proceed by \hoggsay{HOW?}

\paragraph{Implementation details:}
\hoggsay{Rounds of redshift updates followed by RHMF updates.}

\end{document}
